\chapter{IMPLEMENTATION AND VALIDATION GUIDANCE}
\label{app:validation}

This appendix gives a staged procedure for implementing and testing the observer
of Chapter~\ref{ch:dmso} and the controller of Chapter~\ref{ch:control}. The
organizing principle is that each stage must have a falsifiable pass criterion
that is checkable without the next stage being correct. A design that is
assembled all at once and then tuned until it balances provides no evidence
about which component is working.

\section{STAGE 0: PARAMETER IDENTIFICATION}

The detuning of the LQR gains reported in Section~\ref{sec:difficulties} is
evidence of parameter error, and the two parameters most likely responsible are
the pendulum inertia $I_p$ and the CG offset $l$. Both can be measured directly
and should be, before any estimator is judged.

\paragraph{Procedure.} Remove the wheels and suspend the chassis from a knife
edge at the wheel axis so that it swings as a physical pendulum. Record the
small-amplitude period $T$ with the onboard gyroscope. Then
\begin{equation}\label{eq:pendid}
T=2\pi\sqrt{\frac{I_{\mathrm{pivot}}}{M_pgl}},
\qquad
I_{\mathrm{pivot}}=I_p+M_pl^2 .
\end{equation}
Measure $l$ independently by balancing the chassis on a knife edge to locate the
CG, then solve \eqref{eq:pendid} for $I_p$. Repeat at several amplitudes to
confirm the small-angle regime.

\paragraph{Motor constants.} Measure $R$ by four-wire resistance at a dozen
shaft positions and average, since commutator position changes the reading.
Obtain $k_e$ by back-driving the motor at a measured rate and reading the
open-circuit terminal voltage; $k_e$ from the slope of $V$ against
$\dot\vartheta$ is far more reliable than $k_e$ inferred from catalog free-run
and stall figures. For a brushed DC motor in SI units $k_m=k_e$, and any
substantial discrepancy between a measured $k_e$ and a catalog-derived $k_m$
indicates that one of the two is wrong.

\paragraph{Deadzone and backlash.} Ramp the commanded voltage slowly in both
directions with the wheels off the ground and record the breakaway voltages
$\delta_+$ and $\delta_-$ of \eqref{eq:deadzone}. Measure backlash by holding
the motor shaft and rotating the wheel through the deadband, reading the
free rotation directly; this gives $\Delta_b$ in \eqref{eq:backlash}.

\rev{\paragraph{Sign conventions.} Hold the wheels on the ground and rotate the
body by hand in the positive sense of $\phi$ (counterclockwise in
Figure~\ref{fig:fbdpend}, center of gravity moving toward $-x$). The logged tilt
should increase, and the encoder count should change in the same sense as when
the robot rolls forward, since the encoders read $x/r+\phi$
(Section~\ref{sec:simmodel}). These two signs fix the tilt compensation of the
encoder position, $x=\mathrm{pos}-r\phi$, used in the estimators and in the
replays of Section~\ref{sec:implresults}; with either reversed, the compensation
doubles the tilt term instead of removing it.}

\paragraph{Pass criterion.} The LQR gains computed from the measured parameters
stabilize the hardware without manual detuning. If they do not, the model is
still wrong and no amount of online estimation will make the comparison in
Chapter~\ref{ch:results} interpretable.

\section{STAGE 1: OFFLINE ESTIMATOR REPLAY}

Before the observer is allowed to influence any control signal, run it open
loop on recorded data.

\begin{enumerate}[label=\arabic*.,leftmargin=*]

\item Drive the robot with the LQR controller alone and log, at the full loop
rate, the raw gyroscope and accelerometer samples, the raw encoder counts, the
commanded voltage, and a timestamp taken from a hardware timer rather than from
a software counter.

\item Replay the log offline through the DMSO. Because $y_k=x_k$ here, the
observer can be evaluated without closing any loop.

\item Verify the theoretical conditions numerically on the replayed data:
compute $\sigmax(F-\Km)$ for the $\Delta t$ actually achieved, compute
$\phi_{\max}=\max_k\norm{\phi(x_k)}$ realized along the recorded trajectory, and
confirm \eqref{eq:simple}. \rev{This is the step most often skipped. The
$\phi_{\max}$ used to select $\Gamma$ is a design assumption, and the trajectory
is free to violate it.}

\end{enumerate}

\paragraph{Pass criterion.} The estimation error settles inside the predicted
ultimate bound $b_e$ of \eqref{eq:bounds}, and $\Gamma\phi_{\max}^2$ computed
from realized data satisfies \eqref{eq:simple}. If the error settles but the
bound is violated, the gains are outside the proven region and the result is
empirical only.

\section{STAGE 2: OBSERVER GAIN SELECTION}

The conditions of Corollary~\ref{cor:theta1} give a constructive selection
order.

\begin{enumerate}[label=\arabic*.,leftmargin=*]

\item \textbf{Place $A$ first.} Choose $\Km$ so that $\sigmax(F-\Km)$ takes a
target value $\bar\sigma$. Note that $\sigmax$, not the eigenvalues, is what
matters; a gain that places the eigenvalues of $A$ well inside the unit circle
can still have $\sigmax(A)>1$ if $A$ is strongly non-normal. Compute the
singular values explicitly. A practical starting point is
$\bar\sigma\in[0.5,0.8]$. \rev{Smaller values reject model error faster, with a
response time that scales as $1/\log(1/\bar\sigma)$, and lower the model-error
floor $\epsilon_N/(1-\bar\sigma)$ of \eqref{eq:limbounds}. Their cost is
measurement noise, which Theorem~\ref{thm:main} does not model: as $\bar\sigma$
falls, $\Km=F-\bar\sigma I$ approaches $F$ and the prediction follows each noisy
measurement more closely. That side of the trade has to be judged from data.}

\item \textbf{Bound $\phi_{\max}$ by construction.} Rather than estimating
$\phi_{\max}$ from a nominal trajectory, choose activation functions that are
bounded by construction, for example a Gaussian radial basis or a
$\tanh$ polynomial basis, so that $\phi_{\max}$ is known a priori and
independent of how far the trajectory wanders.

\item \textbf{Select $\Gamma$ from \eqref{eq:simple}.} With $\bar\sigma$ and
$\phi_{\max}$ fixed,
\[
\Gamma<\frac{1}{\phi_{\max}^2}\min\left\{\frac{1}{2},\,
\frac{1-\bar\sigma^2}{2\bar\sigma^2}\right\}.
\]
Begin at one tenth of this bound and increase. \rev{The speed of uncertainty
identification is set by $\Gamma$, and the conditions permit any sufficiently
small value, so there is no reason to operate near the boundary.}

\item \textbf{Add $\sigma$-modification.} Use \eqref{eq:sigmod} with
$\kappa$ chosen so that $\Gamma\kappa\ll1$; a decade below $\Gamma$ is a
reasonable start. \rev{This is strongly recommended over relying on persistency
of excitation. During regulation about the upright equilibrium the regressor is
nearly constant and \eqref{eq:pe} will not hold, which is exactly the condition
under which weight drift occurs. Drift is insidious: the state estimate remains
good while the weights wander, and the failure appears only when the trajectory
eventually excites a direction the drifted weights misrepresent.}

\end{enumerate}

\paragraph{Diagnostic.} Log $\norm{\What_k}_F$ throughout every run. A weight
norm that grows slowly and monotonically while the estimation error stays small
is the signature of drift, and it is the reason to add leakage rather than to
increase $\Gamma$.

\section{STAGE 3: CONTROLLER BRING-UP}

The controller of Section~\ref{sec:flatcf} \rev{(the structure of
Section~\ref{sec:cffb} in flat-output coordinates)} has four damping gains,
three filter bandwidths, four learning rates, and four leakage gains. Bring them up in this
order, and only in this order.

\begin{enumerate}[label=\arabic*.,leftmargin=*]

\item \textbf{LQR alone.} Confirm stable regulation with the measured
parameters from Stage~0. Record the baseline metrics of
Section~\ref{sec:testmatrix}.

\item \textbf{Backstepping with networks disabled.} Set $\hat F_1=\hat F_2=0$
and run the structure of Section~\ref{sec:flatcf} with $u_e$ active. This tests
the command filters, the error coordinate bookkeeping, and the sign conventions
without any adaptation. \rev{Most implementation errors in a backstepping design
are sign errors in the cross-term cancellations, and they are far easier to find
with the adaptation off.} \rev{Verify numerically that the implemented error
dynamics are the derived ones: difference the logged $z_i$ to estimate
$\dot z_i$, evaluate the right-hand sides of the error equations of
Section~\ref{sec:flatcf} from the logged signals, and confirm that the two agree
to within the discretization error. A sign error in a cancelling pair, such as
$a_2z_3$ in $\dot z_2$ against $-(a_2/w)z_2$ in $\dot z_3$, appears as a residual
proportional to one of the coordinates. (Logging expressions such as
$z_1z_2-z_2z_1$ is not a test: they are zero for any numbers.)} If the two do not
agree, the implementation does not match the derivation.

\item \textbf{Filter bandwidths.} Raise $\varpi_i$ until $\chi_i$ is small
relative to the corresponding $z_i$, then stop. \rev{With noisy measurements
$\chi_i$ stops shrinking, and can grow, as the bandwidth rises
(Figure~\ref{fig:chi}b), so raise $\varpi_i$ only while $\chi_i$ falls.} The upper limit is the lesser of
one fifth of the sample rate and the lowest unmodeled structural mode of the
chassis. \rev{The stiffening rods described in
Section~\ref{sec:difficulties} exist because that mode was originally low enough
to contaminate the inertial measurements; it should be measured, not assumed,
by striking the chassis and recording the gyroscope response.}

\item \textbf{One network at a time.} Enable network one with network two
disabled, then the reverse, then both. Each network has its own error
coordinate driving it, $z_2$ and $z_4$ respectively, so a network that is not
helping shows up as a failure to reduce its own coordinate.

\item \textbf{Leakage last.} Raise $\kappa_{ij}$ only as far as needed to hold
the weight norms flat over a long run. Every increase adds
$\kappa_{ij}W_{ijm}^2/2$ to the residual $\mathcal{C}$ in
\eqref{eq:Vdotfinal} and therefore to the steady-state error.

\end{enumerate}

\section{STAGE 4: HARDWARE ARCHITECTURE}

Two limitations identified in Section~\ref{sec:difficulties} are architectural
rather than algorithmic, and both bound what any estimator can achieve.

\subsection{Separate the control path from the telemetry path}

Serial transmission dominated the loop budget. The remedy is to make the control
path deterministic and the telemetry path best-effort:

\begin{itemize}[leftmargin=*]
\item Run the control law from a hardware timer interrupt at fixed period, and
do no I/O inside it beyond sensor reads and the actuator write.
\item Write telemetry into a ring buffer from the interrupt and drain the buffer
from the main loop, so that a slow or blocked link cannot extend the control
period.
\item Log to onboard storage at full rate and stream a decimated subset, rather
than streaming everything.
\item Timestamp from a free-running hardware counter, not from a software
accumulator, so that jitter is measurable after the fact. \rev{Report the
achieved loop period distribution, not its nominal value. The discretization
\eqref{eq:zoh} assumes a fixed $\Delta t$, and jitter enters as an unmodeled
disturbance that the networks will attempt, and fail, to explain.}
\end{itemize}

\subsection{The encoder placement limits the achievable result}

The encoders are on the motor shaft, ahead of the gearbox, so backlash is
structurally unobservable. \rev{Increasing the loop rate does not help: the
measurement is not noisy, it is measuring the wrong thing.} Two remedies are
available, in increasing order of effort:

\begin{enumerate}[label=(\alph*),leftmargin=*]
\item \textbf{Load-side encoding.} Add a second encoder on the output shaft or
the wheel. The difference between motor-side and load-side angle is then a
direct measurement of the backlash state, which converts an unmodeled
nonlinearity into an observable quantity and permits an explicit backlash
inverse. This is the single highest-value hardware change available.
\item \textbf{Deterministic timing in hardware.} Implementing the fixed-point
control law in programmable logic removes scheduler jitter entirely and permits
loop rates well beyond what the microcontroller supports. This is worth doing
only after (a); a faster loop around a blind measurement does not improve the
estimate.
\end{enumerate}

\subsection{The sample rate trade}

\rev{There is an interior optimum in $\Delta t$ and it should be found rather
than assumed. Shortening $\Delta t$ improves the fidelity of \eqref{eq:zoh} and
reduces the interval over which the zero-order hold is in error, but the encoder
velocity quantization in Table~\ref{tab:noise} scales as
$1.47\times10^{-4}/\Delta t$, so the velocity measurement degrades linearly as
the loop speeds up. Sweep $\Delta t$ in replay over the logged data from
Stage~1, plot the RMS velocity estimation error against $\Delta t$, and operate
at the minimum. If the minimum lies at the shortest $\Delta t$ tested, the
encoder resolution is not the binding constraint and the loop rate may be raised
further; if it lies in the interior, raising the loop rate requires raising the
encoder resolution first.}

\section{STAGE 5: FALSIFICATION TESTS}

Each of the following is designed to fail if a specific claim is untrue.

\begin{enumerate}[label=T\arabic*.,leftmargin=*]

\item \textbf{Bound verification.} Run case O2 with a deliberately inflated
parameter error and confirm that the estimation error settles inside $b_e$ from
\eqref{eq:bounds} computed with the realized $\phi_{\max}$. A violation means
either the conditions are not met or the implementation does not match the
derivation.

\item \textbf{Small-gain admissibility.} Run the observer at
$\Gamma$ one and two decades below the bound in \eqref{eq:simple}. \rev{Stability
must be retained, with slower identification. This test directly falsifies the
revision~7 condition, which excluded small gains, and confirms the corrected
one.}

\item \textbf{Position drift.} Run cases C2 and C4 for at least ten times the
dominant closed-loop time constant with the two-step and command-filtered
designs side by side, and plot $z_1$ over the full run. \rev{Remark~\ref{rem:twostep} predicts
that the two-step design permits a slow position ramp. Either it appears, which
confirms the analysis, or it does not, which means the LQR term is supplying the
missing damping and should be credited explicitly in the design rather than
absorbed into the approximation targets.}

\item \textbf{Algebraic loop.} Disable the command filters, replacing $x_3^c$
with $\alpha_2$ directly, and attempt to execute the ordering of
Remark~\ref{rem:order}. \rev{The computation should fail to be well posed. This
confirms that the filters are structural rather than cosmetic.}

\item \textbf{Weight drift.} Run a long regulation case with leakage disabled
and log $\norm{\What_k}_F$. Growth in the weight norm with a flat estimation
error confirms that persistency of excitation does not hold in regulation and
justifies $\sigma$-modification.

\item \textbf{Nonsmoothness floor.} Run case C4 with the backlash amplitude
$\Delta_b$ swept over a range including zero. \rev{The residual limit cycle
amplitude should scale with $\Delta_b$ and vanish as $\Delta_b\to0$. If it does
not, the limit cycle has another cause and the explanation in Remark~\ref{rem:backlash} is wrong.}

\end{enumerate}

\section{SOFTWARE ORGANIZATION}

A structure that supports the above is worth stating explicitly.

\begin{itemize}[leftmargin=*]
\item \textbf{One model, three consumers.} The continuous model
\eqref{eq:nl1}--\eqref{eq:nl2}, the discretization \eqref{eq:zoh}, and the
parameter set should exist in exactly one place and be consumed by the truth
propagator, the estimator, and the controller. Divergence between a simulation
model and a flight model is a classic and silent source of error.
\item \textbf{Same source, both targets.} The estimator and controller should
compile for both the desktop simulation and the embedded target from the same
source, with the platform differences confined to sensor and actuator
interfaces. This makes replay meaningful: the code that is tested is the code
that runs.
\item \textbf{Fixed-point discipline early.} If the eventual target is
fixed-point, introduce the scaling in simulation and compare against the
floating-point result, rather than discovering the scaling problems on
hardware.
\item \textbf{Record the conditions, not just the results.} Every run should
emit $\sigmax(A)$, realized $\phi_{\max}$, $\Gamma\phi_{\max}^2$, the achieved
loop period distribution, and the weight norms, alongside the state histories.
These are what make a result checkable a year later.
\end{itemize}
